% =============================================================================
% ABSTRACT
% =============================================================================

\chapter*{Abstract}

Rotating machines play a fundamental role in various industrial sectors, with their reliable operation being essential for process efficiency and safety. Early detection of failures, such as cracks or misalignments, is crucial to avoid unplanned shutdowns and reduce maintenance costs. This work proposes the development of a methodology based on numerical simulations for fault diagnosis in rotating machines, using Higher-Order Spectra (HOS) analysis. The methodology is developed in three main stages: (1) numerical modeling of rotating systems using the ROSS tool to simulate dynamic behavior under normal conditions and with specific faults; (2) application of HOS techniques to simulated data to identify characteristic patterns that differentiate normal conditions from modeled faults; and (3) validation of the obtained results through comparison with data available in the literature. The results demonstrate that bispectrum and trispectrum analysis is effective for identifying non-linear signatures characteristic of different types of faults in rotating systems. The proposed methodology offers a simulation-based alternative that can complement or reduce dependence on physical experiments, contributing to the improvement of fault diagnosis techniques in rotating machines.

\textbf{Keywords:} Rotating machines, Fault diagnosis, Higher-order spectra, Numerical simulation, Vibration analysis.

\cleardoublepage
